\documentclass[11pt,a4paper]{article}

\usepackage[utf8]{inputenc}
\usepackage[T1]{fontenc}
\usepackage{lmodern}
\usepackage{amsmath,amssymb,amsfonts}
\usepackage{booktabs}
\usepackage{hyperref}
\usepackage[margin=1in]{geometry}
\usepackage{microtype}
\usepackage{xcolor}

\hypersetup{
  colorlinks=true,
  linkcolor=blue!60!black,
  citecolor=blue!60!black,
  urlcolor=blue!60!black,
}

\title{\textbf{Three Integers and a Mass:\\Deriving the Standard Model Input Set}}
\author{
  James P.\ Galasyn$^{1}$ and Claude Th\'eodore$^{2}$ \\[6pt]
  \small $^{1}$Software engineer and independent researcher \\
  \small $^{2}$Anthropic, San Francisco, CA
}
\date{}

\begin{document}
\maketitle

\begin{abstract}
Paper~1 derived 23 Standard Model parameters from four inputs ($\alpha$, $m_e$, $m_\mu$, $\Lambda_{\text{tube}}$) and three topological integers ($p=2$, $q=1$, $k=3$). We show that three of those four inputs --- $\alpha$, $m_\mu$, and $\Lambda_{\text{tube}}$ --- are themselves derivable from the three topological integers ($p=2$, $q=1$, $k=R/r=3$) and a single measured mass~$m_e$. The fine-structure constant follows from Skilton's integer right triangle $(88, 105, 137)$, whose generators $(11, 4)$ we identify as null worldtube (NWT) quantum numbers: $p_s = 2k + p^2 + q^2 = 11$, $q_s = p^2 = 4$. The muon mass follows from the Koide formula with topology-determined angle $\theta_K = (6\pi+2)/9$. The proton mass --- and hence $\Lambda_{\text{tube}}$ --- follows from the Lenz formula $m_p/m_e = 6\pi^5$, now structurally explained as $2k \times \pi^{p^2+q^2}$. The complete derivation chain from $(2, 1, 3) + m_e$ to all 23~SM parameters requires no free parameters beyond one mass scale.
\end{abstract}

%----------------------------------------------------------------------
\section{Introduction}
\label{sec:intro}

Paper~1~\cite{paper1} demonstrated that 23 Standard Model parameters --- nine charged fermion masses, three gauge couplings, two Higgs sector parameters, four CKM mixing parameters, and five neutrino parameters --- follow from a single geometric object: a photon confined to a $(2,1)$ torus knot. That derivation required four measured inputs ($\alpha$, $m_e$, $m_\mu$, $\Lambda_{\text{tube}}$) and three topological integers ($p=2$, $q=1$, $k=3$).

The natural question is whether the inputs themselves can be derived. If so, how many are truly independent?

The answer is one. Only the electron mass $m_e$ remains irreducible --- a consequence of dimensional analysis, which requires at least one dimensionful input to set the absolute scale. The remaining three inputs are derivable from the topological integers alone:
\begin{enumerate}
  \item $\alpha$ follows from the Skilton formula $\alpha^{-1} = \sqrt{137^2 + \pi^2}$, where 137 is the hypotenuse of the integer right triangle $(88, 105, 137)$ whose generators are NWT quantum numbers.
  \item $m_\mu$ follows from the Koide formula with the topology-determined angle $\theta_K = (6\pi+2)/9$.
  \item $\Lambda_{\text{tube}}$ follows from the Lenz formula $m_p/m_e = 6\pi^5$, which we decompose as $2k \times \pi^{p^2+q^2}$.
\end{enumerate}

The derivation chain is:
\begin{equation}
  (p, q, k) \xrightarrow{\text{Skilton}} \alpha \xrightarrow{\text{Koide}} m_\mu \xrightarrow{\text{Lenz}} m_p \to \Lambda_{\text{tube}} \xrightarrow{\text{Paper\,1}} \text{23 SM parameters}
  \label{eq:chain}
\end{equation}

Why is $m_e$ special? It is the mass of the lightest charged fermion --- the fundamental $(2,1)$ torus knot. All other masses are related to $m_e$ through geometric and topological ratios. But $m_e$ itself sets the overall scale: the actual size in meters of the electron torus. Dimensional analysis guarantees that no purely topological construction can produce a quantity with dimensions of mass. Nature must supply one number. Everything else is geometry.

%----------------------------------------------------------------------
\section{The Three Integers}
\label{sec:integers}

\subsection{The winding numbers $(p, q) = (2, 1)$}

A $(p,q)$ torus knot winds $p$ times around the hole of the torus (toroidal direction) and $q$ times around the tube (poloidal direction). The choice $(2, 1)$ is uniquely determined by three requirements:
\begin{itemize}
  \item \textbf{Spin-$\tfrac{1}{2}$}: The double toroidal winding requires $4\pi$ rotation to return to the starting point, producing fermionic statistics.
  \item \textbf{Minimal asymmetry}: $p \neq q$ breaks the symmetry between toroidal and poloidal directions, essential for distinguishing electric charge from color charge.
  \item \textbf{Simplest closure}: $q = 1$ is the simplest poloidal winding that closes the curve.
\end{itemize}

The $(2,1)$ torus knot is the unique minimal fermionic knot. No simpler topology produces spin-$\tfrac{1}{2}$ particles.

\subsection{The aspect ratio $k = R/r = 3$}

The third integer is the torus aspect ratio $k = R/r$, where $R$ is the major radius and $r$ the minor radius. In Paper~1, this appeared as $N_c = 3$, interpreted as the number of colors. We now reinterpret it as the aspect ratio of the torus itself --- the physical origin of why $N_c = 3$ in QCD.

The aspect ratio classifies particles by internal structure:

\begin{table}[h]
\centering
\begin{tabular}{@{}clcc@{}}
\toprule
$k$ & Particle type & Pythagorean triple? & Stable? \\
\midrule
1 & leptons & none at $p=1$ & topological \\
2 & mesons & none at $p=1$ & \textbf{no} \\
3 & baryons & \textbf{(3,4,5)} & \textbf{yes} \\
4 & tetraquarks & (4,3,5) & no ($\to$ 2 mesons) \\
5 & pentaquarks & (5,12,13) & resonance \\
\bottomrule
\end{tabular}
\end{table}

The proton is the unique stable baryon because $k=3$ is the lowest aspect ratio admitting a Pythagorean resonance --- the $(3,4,5)$ triple --- at the fundamental winding $p=1$. The resonance condition $(kp)^2 + q^2 = N^2$ at $k=3$, $p=1$, $q=4$ gives $9 + 16 = 25 = 5^2$, a perfect standing wave.

\subsection{Quarks as harmonic modes}

The reinterpretation of $k$ as aspect ratio rather than linking number changes the picture of quarks. In a $k=3$ cavity, the natural harmonic modes have frequencies at $f/k$, $f/k^2$, $f/k^3$ --- that is, $f/3$, $f/9$, $f/27$. Three quarks are three harmonic modes of a single photon cavity, not three independently linked torus knots. The ``color'' of a quark is its harmonic number. The aspect ratio $k=3$ naturally produces three flavors of confined mode --- this is why $N_c = 3$ in QCD.

%----------------------------------------------------------------------
\section{Deriving $\alpha$ from $(p, q, k)$}
\label{sec:alpha}

\subsection{Skilton's formula}

F.~Ray Skilton, in three papers published in the proceedings of the Annual Pittsburgh Conference on Modeling and Simulation (1986--1988)~\cite{skilton1986,skilton1987,skilton1988}, established an integer-based framework for the fundamental constants. His central result:
\begin{equation}
  \alpha^{-1} = \sqrt{137^2 + \pi^2} = 137.036016
  \label{eq:skilton}
\end{equation}

CODATA 2018: $\alpha^{-1} = 137.035999$. Residual: 0.000017. Agreement: \textbf{0.12~ppm} --- accurate to 1~part in 8~million with zero free parameters.

\subsection{The integer right triangle $(88, 105, 137)$}

Skilton identified 137 as the hypotenuse of a Pythagorean triple:
\begin{equation}
  88^2 + 105^2 = 7744 + 11025 = 18769 = 137^2
\end{equation}

The generators of any primitive Pythagorean triple $(a, b, c)$ satisfy:
\begin{equation}
  a = p_s^2 - q_s^2, \quad b = 2p_s q_s, \quad c = p_s^2 + q_s^2
\end{equation}

For $(88, 105, 137)$: the generators are $(p_s, q_s) = (11, 4)$, since $b = 2 p_s q_s = 88$, $a = p_s^2 - q_s^2 = 105$, and $c = p_s^2 + q_s^2 = 137$.

\subsection{The NWT identification}

We identify these generators as NWT quantum numbers:
\begin{align}
  p_s &= 2k + p^2 + q^2 = 2(3) + 4 + 1 = 11, \label{eq:ps} \\
  q_s &= p^2 = 4. \label{eq:qs}
\end{align}

The generator $p_s$ combines the aspect ratio contribution $2k = 6$ with the knot metric $p^2 + q^2 = 5$ to give~11. The generator $q_s$ is simply the toroidal mode count $p^2 = 4$.

With these identifications:
\begin{itemize}
  \item $c = p_s^2 + q_s^2 = 121 + 16 = \mathbf{137}$ (the integer part of $\alpha^{-1}$)
  \item $a = p_s^2 - q_s^2 = 121 - 16 = \mathbf{105}$
  \item $b = 2\,p_s\,q_s = 2 \times 11 \times 4 = \mathbf{88}$
\end{itemize}

\subsection{The $\pi^2$ correction}

The Skilton formula adds $\pi^2$ under the square root:
\begin{equation}
  \alpha^{-1} = \sqrt{c^2 + \pi^2} = \sqrt{137^2 + \pi^2}
\end{equation}

The $\pi^2$ correction arises from the poloidal winding: $q = 1$ contributes one factor of $\pi$ per toroidal transit (the circumference of the tube cross-section), and $p = 2$ toroidal transits give two such contributions. The correction is $\pi^2$ because the photon traverses two poloidal half-circuits, each contributing a phase of~$\pi$. The squared correction reflects the energy, which goes as phase$^2$.

\subsection{Verification}

\begin{table}[h]
\centering
\begin{tabular}{@{}ll@{}}
\toprule
Quantity & Value \\
\midrule
$\alpha^{-1}$ (Skilton) & 137.036016 \\
$\alpha^{-1}$ (CODATA 2018) & 137.035999 \\
Residual & 0.000017 \\
Relative error & $1.2 \times 10^{-7}$ (0.12~ppm) \\
\bottomrule
\end{tabular}
\end{table}

This is the first structural derivation of the fine-structure constant from topology. Previous numerological relationships (Eddington's 136~\cite{eddington1929}, Wyler's formula~\cite{wyler1969}) required ad hoc constructions. Here, 137 emerges as $c = p_s^2 + q_s^2$ from the generators of the unique Pythagorean triple associated with the $(2,1)$ knot on a $k=3$ torus.

\subsection{Uniqueness}

In Part~3 (1988), Skilton performed an exhaustive computer search over all reduced integer right triangles with generators $p, q \leq 100$ and found exactly 29 solutions. He proved that $(88, 105, 137)$ is unique in possessing the number-theoretic properties required to encode the fundamental mass ratios. Our identification of the generators as NWT quantum numbers explains this uniqueness: there is only one minimal fermionic knot ($p=2$, $q=1$) and only one lowest stable baryon aspect ratio ($k=3$), so there is only one triple.

%----------------------------------------------------------------------
\section{Deriving $m_\mu$ from $(p, q, k) + m_e$}
\label{sec:muon}

\subsection{The Koide formula}

The Koide formula~\cite{koide1983} relates the three charged lepton masses:
\begin{equation}
  Q = \frac{m_e + m_\mu + m_\tau}{(\sqrt{m_e} + \sqrt{m_\mu} + \sqrt{m_\tau})^2} = \frac{2}{3}
  \label{eq:koide}
\end{equation}

This holds to 0.0009\% with PDG masses~\cite{pdg2022}. The parametric form is:
\begin{equation}
  \sqrt{m_i} = \frac{S}{3}\left(1 + \sqrt{2}\cos\!\left(\theta_K + \frac{2\pi i}{3}\right)\right)
  \label{eq:koide-param}
\end{equation}
where $S$ is the mass scale and $\theta_K$ is the Koide angle. $Q = 2/3$ is automatic in this parameterization; all the physical content is in~$\theta_K$.

\subsection{The topology-determined angle}

Paper~1 derived the Koide angle from the torus knot topology:
\begin{equation}
  \theta_K = \frac{p}{k}\!\left(\pi + \frac{q}{k}\right) = \frac{6\pi + 2}{9} = 2.31662~\text{rad}
  \label{eq:koide-angle}
\end{equation}

With $k = 3$, the decomposition is:
\begin{itemize}
  \item $2\pi/3$: the $\mathbb{Z}_3$ base symmetry from the $k=3$ aspect ratio (three harmonic modes)
  \item $2/9 = pq/k^2$: the toroidal-poloidal coupling through the $k \times k$ interaction matrix
\end{itemize}

\subsection{$m_\mu$ as derived quantity}

Given $m_e = 0.51100$~MeV and $\theta_K = (6\pi+2)/9$, the Koide parameterization with $Q = 2/3$ determines $m_\mu$ uniquely. The scale parameter $S$ follows from $m_e$ via:
\begin{equation}
  \sqrt{m_e} = \frac{S}{3}\!\left(1 + \sqrt{2}\cos\theta_K\right)
\end{equation}

Then:
\begin{equation}
  \sqrt{m_\mu} = \frac{S}{3}\!\left(1 + \sqrt{2}\cos\!\left(\theta_K + \frac{2\pi}{3}\right)\right)
\end{equation}

\textbf{Result}: $m_\mu = 105.658$~MeV (measured: 105.658~MeV, 0.001\% error).

Paper~1 treated $m_\mu$ as an input only because $\theta_K$ appeared to be an empirical fit. Now that $\theta_K = (6\pi+2)/9$ is derived from the topology $(p, q, k)$, the muon mass follows from $m_e$ alone. It was always ``derivable'' --- we simply did not recognize the derivation until the topology was understood.

%----------------------------------------------------------------------
\section{Deriving $m_p$ and $\Lambda_{\text{tube}}$ from $(p, q, k) + m_e$}
\label{sec:proton}

\subsection{The Lenz formula}

Friedrich Lenz, in a 1951 letter to \emph{Physical Review}~\cite{lenz1951}, noted:
\begin{equation}
  \frac{m_p}{m_e} \approx 6\pi^5 = 1836.118
  \label{eq:lenz}
\end{equation}

Measured: $m_p/m_e = 1836.153$. Error: \textbf{0.002\%}. Like Skilton's $\alpha$ formula, this has zero free parameters and sub-percent accuracy.

\subsection{The NWT structural decomposition}

We identify the structural origin of each factor:
\begin{equation}
  6\pi^5 = 2k \times \pi^{p^2 + q^2} = 2(3) \times \pi^{4+1} = 6\pi^5
  \label{eq:lenz-decomp}
\end{equation}

\begin{table}[h]
\centering
\begin{tabular}{@{}lll@{}}
\toprule
Factor & NWT origin & Value \\
\midrule
2 & $p = 2$ (toroidal winding number) & 2 \\
$k = 3$ & aspect ratio (baryon number) & 3 \\
$2k = 6$ & harmonic mode count: $k$ modes $\times$ 2 spin states & 6 \\
$p^2 + q^2 = 5$ & knot metric of $(2,1)$ torus knot & 5 \\
$\pi^5$ & five-dimensional phase space of $(2,1)$ knot on $k\!=\!3$ torus & 306.0 \\
\bottomrule
\end{tabular}
\end{table}

\subsection{Physical interpretation}

The proton is a $k=3$ cavity containing a confined photon in a $(2,1)$ knot. The factor $2k = 6$ counts the number of distinct harmonic modes: three quarks (the $f/3$, $f/9$, $f/27$ harmonics of the $k=3$ cavity) each with two spin orientations. The $\pi^5$ factor encodes the five-dimensional phase space: the $(2,1)$ knot has two toroidal and one poloidal degree of freedom on the torus surface, plus two additional degrees from the knot metric (the path length $\sqrt{p^2 + q^2}$ introduces coupling between directions). Each degree of freedom contributes a factor of $\pi$ from the poloidal circumference integration.

\subsection{Proton charge radius}

The proton mass determines the proton charge radius through:
\begin{equation}
  m_p R_p = p^2 \hbar c = 4\hbar c
  \label{eq:charge-radius}
\end{equation}

Using the muonic hydrogen proton radius $R_p = 0.841$~fm (post-puzzle value): $m_p \times R_p = 938.3 \times 0.841 = 789.1$~MeV$\cdot$fm versus $p^2 \times \hbar c = 4 \times 197.327 = 789.3$~MeV$\cdot$fm. Error: \textbf{0.03\%}. The winding number squared $p^2 = 4$ sets the charge radius scaling --- the charge distribution extends over $p^2$ Compton wavelengths.

\subsection{From $m_p$ to $\Lambda_{\text{tube}}$}

With $m_p$ determined, the proton torus radius follows from the self-consistency condition:
\begin{equation}
  R_p = \frac{\hbar c}{m_p c^2} \times f(p,q,\alpha) \approx 0.106~\text{fm}
\end{equation}
where $f(p,q,\alpha)$ encodes the ratio of circulation energy to total energy for the $(2,1)$ knot. The tube energy scale is then:
\begin{equation}
  \Lambda_{\text{tube}} = \frac{\hbar c}{\alpha R_p} = 255.7~\text{GeV}
  \label{eq:tube-scale}
\end{equation}

This is the same value used in Paper~1, but now derived rather than measured.

\subsection{Independent route: Cornell potential}

Paper~1 (Section~VII.D--E and Supplementary~S17) provides an independent derivation of $m_p$ via the Cornell potential:
\begin{equation}
  m_p = 2\sqrt{\sigma\,\hbar c\!\left(\frac{3}{2} - 2\alpha_s\right)} = 945.7~\text{MeV} \quad (0.8\%~\text{error})
\end{equation}
where $\sigma = k^2 m_\pi^2/\hbar c = 894.6$~MeV/fm and $\alpha_s = 16\alpha$. This route uses two variational coefficients ($3/2$ and~2) not yet derived from NWT dynamics. The Lenz route is more direct and more accurate (0.002\% vs 0.8\%), but the Cornell route has the advantage of illuminating the energy budget (kinetic 59\%, Coulomb $-9$\%, string 50\%).

%----------------------------------------------------------------------
\section{The Complete Derivation Chain}
\label{sec:chain}

The full chain from $(p, q, k) + m_e$ to all 23 Standard Model parameters:
\begin{equation}
  (2, 1, 3) + m_e \xrightarrow{\text{Skilton}} \alpha \xrightarrow{\text{Koide}} m_\mu \xrightarrow{\text{Lenz}} m_p \to \Lambda_{\text{tube}} \xrightarrow{\text{Paper\,1}} \text{23 SM parameters}
\end{equation}

\begin{table}[h]
\centering
\caption{Derivation of former inputs from $(p, q, k) = (2, 1, 3)$.}
\label{tab:derivation}
\begin{tabular}{@{}llllc@{}}
\toprule
Former input & Formula & Predicted & Measured & Error \\
\midrule
$\alpha^{-1}$ & $\sqrt{(p_s^2+q_s^2)^2+\pi^2}$ & 137.036016 & 137.035999 & 0.12~ppm \\
$m_\mu$ & Koide($\theta_K$, $m_e$) & 105.658~MeV & 105.658~MeV & 0.001\% \\
$m_p/m_e$ & $2k \times \pi^{p^2+q^2}$ & 1836.118 & 1836.153 & 0.002\% \\
$\Lambda_{\text{tube}}$ & $\hbar c/(\alpha R_p)$ & 255.7~GeV & 255.7~GeV & derived \\
\bottomrule
\end{tabular}
\end{table}

\begin{table}[h]
\centering
\caption{Complete parameter count.}
\label{tab:count}
\begin{tabular}{@{}lr@{}}
\toprule
Description & Count \\
\midrule
\textbf{Irreducible inputs} & \\
\quad Measured mass: $m_e$ & 1 \\
\quad Topological integers: $p$, $q$, $k$ & 3 \\
\textbf{Total inputs} & \textbf{4} \\
\midrule
\textbf{Derived quantities} & \\
\quad Former inputs now derived: $\alpha$, $m_\mu$, $\Lambda_{\text{tube}}$ & 3 \\
\quad SM parameters (Paper~1): masses, couplings, mixing & 23 \\
\textbf{Total predictions} & \textbf{26} \\
\bottomrule
\end{tabular}
\end{table}

The Standard Model's 26 parameters reduce to $(p, q, k, m_e) = (2, 1, 3, 0.511~\text{MeV})$.

%----------------------------------------------------------------------
\section{Geodesic Curvature and the Junction Condition}
\label{sec:junction}

\subsection{The photon is never a geodesic}

The $(2,1)$ torus knot has geodesic curvature:
\begin{equation}
  \kappa_g \times R = \frac{1}{\sqrt{2}}
\end{equation}
on a torus with aspect ratio $r/R = \alpha$. A true geodesic on the torus surface would have $\kappa_g = 0$. The photon path has nonzero geodesic curvature everywhere --- it is \emph{not} a geodesic of the torus surface. Rather, it is a null geodesic of the ambient spacetime that is \emph{confined} to the torus surface by electromagnetic self-interaction.

\subsection{Israel junction conditions}

The Israel junction conditions~\cite{israel1966} for a null shell on the torus surface reduce to the Young-Laplace equation~\cite{young1805,laplace1805} --- a balance between electromagnetic radiation pressure and surface tension:
\begin{equation}
  [K_{ab}] = 8\pi G\, S_{ab}
\end{equation}
where $[K_{ab}]$ is the jump in extrinsic curvature across the shell and $S_{ab}$ is the surface stress-energy. For the self-consistent NWT torus, this reduces to:
\begin{equation}
  \Delta P = \frac{\sigma_{\text{surface}}}{R_{\text{curv}}}
\end{equation}
--- the same equation governing soap bubbles and fluid interfaces.

\subsection{The self-energy equation IS the junction condition}

A crucial finding: the Israel junction condition does not provide an independent equation beyond the electromagnetic self-energy balance already used in Paper~1 (Section~II). The self-consistency condition $r/R = \alpha$ already encodes the force balance. There is no ``second equation'' constraining the aspect ratio from the junction conditions alone.

This means the aspect ratio $k = R/r$ must be fixed by the topology --- specifically, by which harmonic modes can exist as standing waves on the torus surface. The junction condition determines the \emph{shape} ($r/R = \alpha$ for the electromagnetic aspect ratio), while the \emph{harmonic structure} ($k = 3$ for baryons) is determined by the Pythagorean resonance condition. These are independent constraints operating at different scales.

%----------------------------------------------------------------------
\section{What Remains}
\label{sec:open}

\subsection{$m_e$ is irreducible}

The electron mass $m_e = 0.511$~MeV is the one number that cannot be derived from topology. Dimensional analysis requires at least one dimensionful constant to set the absolute scale of the universe. Whether $m_e$ can be derived from $c$, $\hbar$, and $G$ (which themselves define the unit system) is beyond the scope of this framework.

\subsection{Open problems}

Several links in the derivation chain require further theoretical development:
\begin{enumerate}
  \item \textbf{The $\pi^2$ correction}: Skilton's formula $\alpha^{-1} = \sqrt{137^2 + \pi^2}$ works to 0.12~ppm, and we have identified the generators as NWT quantum numbers. But the origin of the $\pi^2$ correction --- why $\alpha^{-1}$ involves a \emph{quadrature sum} of the integer hypotenuse and~$\pi$ --- needs a deeper geometric derivation from the torus mode structure.

  \item \textbf{The Koide angle derivation}: $\theta_K = (6\pi+2)/9$ is derived from $(p, q, k)$ through a specific formula (Paper~1, Section~III.A), but a complete spectral-geometry proof --- showing that this angle is the \emph{unique} solution of the eigenvalue problem on the $(2,1)$ torus --- is not yet available.

  \item \textbf{The Lenz decomposition}: $6\pi^5 = 2k \times \pi^{p^2+q^2}$ is verified numerically and the physical interpretation (harmonic modes $\times$ phase space) is compelling, but a rigorous derivation from the torus wave equation --- showing that the proton-to-electron mass ratio must equal this product --- requires more work.

  \item \textbf{The 0.12~ppm residual}: The Skilton $\alpha$ formula has an absolute residual of 0.000017 (0.12~ppm) from CODATA. This may be a higher-order correction from radiative effects (the Lamb shift, anomalous magnetic moment, \emph{etc.})\ or may indicate that the formula is approximate. Whether the residual is structurally meaningful or coincidental is unknown.

  \item \textbf{Quark mass hierarchy}: The quark Koide extension uses a hierarchy parameter $B^2$ that is approximate at the ${\sim}3\%$ level. Deriving $B^2$ from $(p, q, k)$ with the same precision as the lepton sector remains open.
\end{enumerate}

%----------------------------------------------------------------------
\section{Conclusion}
\label{sec:conclusion}

The Standard Model's 26 free parameters reduce to three integers and one mass:
\begin{equation}
  (p, q, k, m_e) = (2, 1, 3, 0.511~\text{MeV})
\end{equation}

The $(2,1)$ torus knot is the simplest topology producing spin-$\tfrac{1}{2}$ fermions. The aspect ratio $k = 3$ is the lowest value admitting a Pythagorean resonance $(3,4,5)$, making the proton the unique stable baryon. The electron mass $m_e$ is the one number nature must supply --- it sets the absolute scale. Everything else is geometry:
\begin{itemize}
  \item $\alpha$ from the Skilton formula with NWT generators (0.12~ppm)
  \item $m_\mu$ from the Koide formula with topology-determined angle (0.001\%)
  \item $m_p$ from the Lenz formula with NWT decomposition (0.002\%)
  \item 23 SM parameters from Paper~1's derivation chain (median 0.7\%)
\end{itemize}

The question ``why these 26 parameters?'' reduces to ``why this knot?'' --- and the answer is: $(2,1)$ is the simplest fermion, $k=3$ is the first stable baryon, and $m_e$ is the scale of the universe. There is nothing else to explain.

%----------------------------------------------------------------------
\begin{thebibliography}{12}

\bibitem{paper1}
J.~P.~Galasyn and C.~Th\'eodore,
``The Standard Model from a Torus Knot: Spectrum, Resonance Structure, and Decay Dynamics,''
Paper~1 in this series (2026).

\bibitem{pdg2022}
Particle Data Group, R.~L.~Workman \emph{et al.},
``Review of particle physics,''
\emph{Prog.\ Theor.\ Exp.\ Phys.}\ \textbf{2022}, 083C01 (2022); 2024 update.

\bibitem{skilton1986}
F.~R.~Skilton,
``Foundation for an integer-based cosmological model,''
\emph{Proc.\ 17th Annual Pittsburgh Conf.\ on Modeling and Simulation},
Vol.~17, Part~1, pp.~295--300 (1986).

\bibitem{skilton1987}
F.~R.~Skilton,
``Foundation for an integer-based cosmological model, Part~2 --- Evenness,''
\emph{Proc.\ 18th Annual Pittsburgh Conf.\ on Modeling and Simulation},
Vol.~18, Part~5, pp.~1623--1630 (1987).

\bibitem{skilton1988}
F.~R.~Skilton,
``Foundation for an integer-based cosmological model, Part~3 --- Integers and the natural constants,''
\emph{Proc.\ 19th Annual Pittsburgh Conf.\ on Modeling and Simulation},
Vol.~19, Part~1, pp.~9--12 (1988).

\bibitem{koide1983}
Y.~Koide,
``New viewpoint on quark and lepton mass hierarchy,''
\emph{Phys.\ Rev.\ D}\ \textbf{28}, 252 (1983).

\bibitem{lenz1951}
F.~Lenz,
``The ratio of proton and electron masses,''
\emph{Phys.\ Rev.}\ \textbf{82}, 554 (1951).

\bibitem{eddington1929}
A.~S.~Eddington,
``The charge of an electron,''
\emph{Proc.\ R.\ Soc.\ A}\ \textbf{122}, 358 (1929).

\bibitem{wyler1969}
A.~Wyler,
``L'espace sym\'etrique du groupe des \'equations de Maxwell,''
\emph{C.\ R.\ Acad.\ Sci.\ Paris}\ \textbf{269}, 743 (1969).

\bibitem{israel1966}
W.~Israel,
``Singular hypersurfaces and thin shells in general relativity,''
\emph{Nuovo Cimento B}\ \textbf{44}, 1 (1966).

\bibitem{young1805}
T.~Young,
``An essay on the cohesion of fluids,''
\emph{Phil.\ Trans.\ R.\ Soc.\ Lond.}\ \textbf{95}, 65 (1805).

\bibitem{laplace1805}
P.~S.~Laplace,
``Suppl\'ement au dixi\`eme livre du Trait\'e de m\'ecanique c\'eleste,''
in \emph{Trait\'e de M\'ecanique C\'eleste}, Vol.~4 (Courcier, 1805).

\end{thebibliography}

\end{document}
